\section{Conclusion}\label{sec:conclusion}
In the signed-Radon framework of Li and Zou \cite{LiZou},
Theorem~\ref{thm:finite} gives stability of extended minimum values for
convergent weights under eventual relative Gram comparison with factors
tending to one. At finite
target cost, the minimum is attained, near-minimizers admit optimal weak-star
clusters, and Frobenius variations vanish in the null directions of the
limiting weights. The fixed-basis, reaction-zero and partial-reaction results
are specializations of this comparison mechanism.

For large transport scalings with fixed reaction weight, the limiting
reaction-channel infimum retains the original finite-baseline class.
Theorems~\ref{thm:ray} and~\ref{thm:uniform} treat fixed, possibly singular
rays and uniform divergence in every direction separately. For entrywise
divergence on a convex domain with positive-definite reaction weight, the hub
construction identifies the latter limit with the weighted Bures value of
the total endpoint matrices and gives an explicit recovery bound.

The rotating families and norm-divergent sequence exhibit failures of value
continuity and of magnitude-based convergence assumptions. The unscaled-flux
example separates convergence of actions from compactness of fluxes. The
Kantorovich--Bures obstruction and smooth positive separation show that the
operator can affect both large-weight limits and costs along a common
weak-star convergent family of uniformly positive densities.
